\clearpage
\pdfbookmark[2]{3. Bestimmung der Elemente}{parabolisch-003}
\begin{center}
[3.]\par
\medskip
Zur parabolischen Bewegung.\par
\medskip
[Bestimmung der Elemente durch zwei Radios Vectores \(r\), \(r'\) und Chorde \(\rho\).]
\end{center}

Es sei
\[
\mathrm{[1]}\qquad \frac{r'-r}{\rho}=\sin\psi,\qquad \frac{\rho}{r'+r}=\sin\varphi.
\]
\SourceLeaf{334}[Mit den Bezeichnungen der Theoria motus Art. 98 und 108 hat man:
\[
\begin{gathered}
\frac{p(r'+r)}{2rr'}=1+\cos f\cos F,\qquad
\frac{p(r'-r)}{2rr'}=\sin f\sin F\\
\frac{pp}{rr'}=(\cos f+\cos F)^2;
\end{gathered}
\]
also
\[
\begin{gathered}
\frac{r'+r}{2p}=\frac{1+\cos f\cos F}{(\cos f+\cos F)^2}\mathpunct{\raisebox{.45ex}{,}}\\
\frac{r'-r}{2p}=\frac{\sin f\sin F}{(\cos f+\cos F)^2}
\end{gathered}
\]
oder, da
\[
p=\frac{\rho}{\eta(\theta'-\theta)}\mathpunct{\raisebox{.45ex}{,}}
\]
wo
\[
\begin{gathered}
\eta\eta=1+\tfrac14(\theta'+\theta)^2=\frac{1+\cos f^2+2\cos f\cos F}{(\cos f+\cos F)^2}\mathpunct{\raisebox{.45ex}{,}}\\
\theta'-\theta=\frac{2\sin f}{\cos f+\cos F}\mathpunct{\raisebox{.45ex}{,}}
\end{gathered}
\]
so wird
\[
\begin{gathered}
\frac{\rho}{r'+r}=\sin\varphi=\frac{\sin f\BareRoot{(1+\cos f^2+2\cos f\cos F)}}{1+\cos f\cos F}\\
\frac{r'-r}{\rho}=\sin\psi=\frac{\sin F}{\BareRoot{(1+\cos f^2+2\cos f\cos F)}}\mathpunct{\cdot}
\end{gathered}
\]

Hieraus:
\[
\begin{gathered}
\cos\varphi=\frac{\cos f(\cos f+\cos F)}{1+\cos f\cos F}\mathpunct{\raisebox{.45ex}{,}}\\
\cos\psi=\frac{\cos f+\cos F}{\BareRoot{(1+\cos f^2+2\cos f\cos F)}}\mathpunct{\raisebox{.45ex}{,}}
\end{gathered}
\]
also]
\[
\mathrm{[2]}\qquad \cos\psi\,\mathrm{tang}\,\varphi=\mathrm{tang}\,f
\]
[woraus \(f\) gefunden wird.

Ferner
\[
\begin{gathered}
\sin\varphi\sin\psi=\frac{\sin f\sin F}{1+\cos f\cos F}\\
\sin\varphi\cos\psi=\frac{\sin f(\cos f+\cos F)}{1+\cos f\cos F}\mathpunct{\raisebox{.25ex}{;}}
\end{gathered}
\]
hieraus
\[
\begin{gathered}
\sin\varphi\sin(F-\psi)=\cos f\sin\varphi\sin\psi=\sin f\,\mathrm{tang}\,\psi\cos\varphi\\
\sin\varphi\cos(F-\psi)=\sin f.
\end{gathered}
\]

Diese Gleichungen selbst, oder die daraus folgenden:]
\[
\mathrm{tang}\,\psi\cos\varphi=\mathrm{tang}\,(F-\psi)
\]
oder
\[
\begin{gathered}
\sin f\sin\psi\,\mathrm{tang}\,\tfrac12\varphi=\sin(2\psi-F)\\
\sin f\sin\psi\,\mathrm{cotang}\,\tfrac12\varphi=\sin F
\end{gathered}
\]
\SourceLeaf{335}[dienen zur Bestimmung von \(F\);] die wahren Anomalien [sind dann] \(F-f\) und \(F+f\).

[Weiter ist
\[
\frac{r'+r}{2p}=\frac{\sin f}{\sin F}\frac{\mathrm{tang}\,\psi}{\sin\varphi\cos\psi}=\frac{\mathrm{tang}\,\tfrac12\varphi}{\sin\varphi\cos\psi^2}\mathpunct{\raisebox{.45ex}{,}}
\]
also]
\[
q=\left[\frac p2=\right]\tfrac12(r'+r)\cos\tfrac12\varphi^2\,.\cos\psi^2,
\]
[woraus \(q\) bekannt ist.

Zur Berechnung der seit dem Periheldurchgange verflossenen Zeiten hat man mit der Bezeichnung der Notiz 1:
\[
\begin{aligned}
mm&=2(r'+r)\cos(45^0-\tfrac12\varphi)^2\\
nn&=2(r'+r)\sin(45^0-\tfrac12\varphi)^2\\
mn&=(r'+r)\cos\varphi
\end{aligned}
\]
also
\[
\begin{aligned}
m+n&=2\BareRoot{(r'+r)}\,.\cos\tfrac12\varphi\\
m-n&=2\BareRoot{(r'+r)}\,.\sin\tfrac12\varphi.
\end{aligned}
\]
Da nun
\[
r'-r=(r'+r)\sin\psi\sin\varphi,
\]
so wird]
\[
\tfrac12(t'+t)=\frac{(r'+r)^{\frac32}\sin\psi\cos\tfrac12\varphi}{6k}(3-2\sin\psi^2\cos\tfrac12\varphi^2).
\]
\begin{center}\rule{.3\linewidth}{.4pt}\end{center}
